\documentclass[a4paper,11pt]{article}
\usepackage[T1]{fontenc}
\usepackage[utf8]{inputenc}
\usepackage[italian]{babel}
\usepackage{amsmath,amssymb}
\usepackage{graphicx}
\usepackage{geometry}
\usepackage{xcolor}
\usepackage{listings}
\geometry{margin=2cm}

\lstdefinelanguage{MiniZinc}{
  morekeywords={int,set,of,array,var,constraint,include,solve,satisfy,output,show},
  sensitive=true,
  morecomment=[l]{\%},
  morestring=[b]"
}

\lstset{
  language=MiniZinc,
  basicstyle=\ttfamily\small,
  keywordstyle=\color{blue}\bfseries,
  commentstyle=\color{teal!70!black}\itshape,
  stringstyle=\color{orange!70!black},
  numbers=left,
  numberstyle=\tiny\color{gray},
  stepnumber=1,
  numbersep=8pt,
  showstringspaces=false,
  frame=single,
  breaklines=true
}

\title{CSP: Generalised Arc-Consistency}
\date{}

\begin{document}
\maketitle

\subsection*{1.1 Pseudocodice: nodo-consistenza}

\begin{lstlisting}
def NODO-CONSISTENZA(CSP): ritorna falso se trova inconsistenza
  per ogni variabile X in CSP:
    per ogni valore a in D(X):
      se a non soddisfa tutti i vincoli unari su X allora
        D(X) = D(X) \ {a}
        se D(X) e' vuoto allora
          ritorna falso
  ritorna vero
\end{lstlisting}


\subsection*{1.2 Pseudocodice: arc-consistenza}
\begin{lstlisting}
def GAC-3(CSP): ritorna falso se trova inconsistenza
  Q = {(Xi, c) | c e' un vincolo in CSP che coinvolge Xi}
  while Q non e' vuoto:
    (Xi, c) = Q.pop()
    se RIMUOVI_VAL_INCONSISTENTI(Xi, c) allora
      se D(Xi) e' vuoto allora
        ritorna falso
      per ogni vincolo c' in CSP che coinvolge Xi e c' != c:
        per ogni variabile Xj coinvolta in c' diversa da Xi:
          Q.add((Xj, c'))
  ritorna vero

def RIMUOVI_VAL_INCONSISTENTI(Xi, c): ritorna vero se ha rimosso almeno un valore
  modificato = falso
  per ogni valore v in D(Xi):
    se non esiste una tupla di valori per le altre variabili coinvolte in c che soddisfa c quando Xi=v allora
      D(Xi) = D(Xi) \ {v}
      modificato = vero
  ritorna modificato
\end{lstlisting}

\newpage

\subsection*{1.3 Esecuzione su CSP fornito}

Il CSP fornito è il seguente:

\textbf{Variabili:} $X = \{X_1, \dots,  X_5\}$


\textbf{Domini:} $D_1 = \dots = D_5 = \{1, 2, 3, 4, 5\}$


\textbf{Vincoli:} $C = \{C_1, \dots, C_5\}$, dove:

\begin{itemize}
  \item $C_1: X_1 > X_3$
  \item $C_2: X_2 \leq X_3$
  \item $C_3: X_3^2 + X_4^2 \leq 15$
  \item $C_4: X_5 \geq 3$
  \item $C_5: X_1 + X_5 \geq 3$
\end{itemize}

Esecuzione di \texttt{NODO-CONSISTENZA}:

\begin{itemize}
  \item $C_4$ è un vincolo unario su $X_5$, quindi rimuove i valori $1$ e $2$ da $D(X_5)$.
\end{itemize}

Dopo l'esecuzione di \texttt{NODO-CONSISTENZA}, i domini sono:
\[
D(X_1) = D(X_2) = D(X_3) = D
(X_4) = \{1, 2, 3, 4, 5\}, \quad D(X_5) = \{3, 4, 5\}.
\]

Esecuzione di \texttt{GAC-3}:

\begin{itemize}
  \item Q = $\{(X_1, C_1), (X_3, C_1), (X_2, C_2), (X_3, C_2), (X_3, C_3), (X_4, C_3), (X_5, C_4), (X_1, C_5), (X_5, C_5)\}$
  \item Rimuoviamo $(X_1, C_1)$ da Q. Rimuoviamo $1$ da $D(X_1)$. Aggiungiamo $(X_5, C_5)$ a Q (già presente).
  \subitem $D(X_1) = \{2, 3, 4, 5\}$
  \item Rimuoviamo $(X_3, C_1)$ da Q. Rimuoviamo $5$ da $D(X_3)$. Aggiungiamo $(X_2, C_2)$ e $(X_4, C_3)$ a Q (già presenti).
  \subitem $D(X_3) = \{1, 2, 3, 4\}$
  \item Rimuoviamo $(X_2, C_2)$ da Q. Rimuoviamo $5$ da $D(X_2)$. Non aggiungiamo nulla a Q.
  \subitem $D(X_2) = \{1, 2, 3, 4\}$
  \item Rimuoviamo $(X_3, C_2)$ da Q. Non rimuoviamo nulla da $D(X_3)$. Non aggiungiamo nulla a Q.
  \item Rimuoviamo $(X_3, C_3)$ da Q. Rimuoviamo $4$ da $D(X_3)$. Aggiungiamo $(X_1, C_1)$ e $(X_2, C_2)$ a Q.
  \subitem $D(X_3) = \{1, 2, 3\}$, Q = $\{(X_4, C_3), (X_5, C_4), (X_1, C_5), (X_5, C_5), (X_1, C_1), (X_2, C_2)\}$
  \item Rimuoviamo $(X_4, C_3)$ da Q. Rimuoviamo $4$ e $5$ da $D(X_4)$. Non aggiungiamo nulla a Q.
  \subitem $D(X_4) = \{1, 2, 3\}$
  \item Rimuoviamo $(X_5, C_4)$ da Q. Non rimuoviamo nulla da $D(X_5)$. Non aggiungiamo nulla a Q.
  \item Rimuoviamo $(X_1, C_5)$ da Q. Non rimuoviamo nulla da $D(X_1)$. Non aggiungiamo nulla a Q.
  \item Rimuoviamo $(X_5, C_5)$ da Q. Non rimuoviamo nulla da $D(X_5)$. Non aggiungiamo nulla a Q.
  \item Rimuoviamo $(X_1, C_1)$ da Q. Non rimuoviamo nulla da $D(X_1)$. Non aggiungiamo nulla a Q.
  \item Rimuoviamo $(X_2, C_2)$ da Q. Rimuoviamo $4$ da $D(X_2)$. Non aggiungiamo nulla a Q.
  \subitem $D(X_2) = \{1, 2, 3\}$
\end{itemize}

Dopo l'esecuzione di \texttt{GAC-3}, i domini sono:
\[
D(X_1) = \{2, 3, 4, 5\}, \quad D(X_2) = \{1, 2, 3\}, \quad D(X_3) = \{1, 2, 3\}, \quad D(X_4) = \{1, 2, 3\}, \quad D(X_5) = \{3, 4, 5\}.
\]

\subsection*{1.4 Soddisfacibilità}
L'arc-consistenza è una condizione necessaria ma non sufficiente per la soddisfacibilità di un CSP. 
Il CSP fornito è arc-consistente, ma non è detto quindi sia soddisfacibile. Per verificare la soddisfacibilità, è necessario eseguire una ricerca. 
\end{document}